\documentclass[instructions.tex]{subfiles}
\begin{document}
 
\chapter{Des parures et des habillemens immodestes.}
 
On peut dire, sans crainte d'exagérer, que les personnes qui se les permettent s'avilissent elles-mêmes; qu'elles chargent leur conscience de beaucoup de péchés, sans presque même qu'elles s'en doutent, et que rien ne les excusera devant Dieu, lorsqu'elles paroîtront à son redoutable tribunal.
 
\primo En effet, ces excès sont proscrits par l'Écriture, les saints Pères, les païens mêmes; ils sont opposés à l'esprit du christianisme, au bonheur des familles, et ils attirent un juste mépris sur ceux qui s'y abandonnent.
 
Saint Paul veut que la modestie des chrétiens paroisse aux yeux de tous les hommes\footnote{Philipp. 4. 5.}. Il recommande surtout aux personnes du sexe une tenue où brillent la simplicité, la modération et la décence. Il va même jusqu'à leur interdire certains ornemens très à la mode de nos jours, mais qui ne servent qu'à nourrir l'orgueil et à fomenter la vanité. Ecoutons-le parler lui-même : Que les femmes, lorsqu'elles se parent, le fassent avec honnêteté et modestie, sans frisure dans leurs cheveux, sans or, sans perles, sans étoffes de prix; mais selon qu'il convient à des femmes qui veulent donner des marques de piété par de bonnes œuvres\footnote{Tim. 2. 9, 10.}. Le Prince des apôtres ne montre pas, sur ce point, moins de sollicitude et de zèle que le Docteur des nations. Il invite les femmes chrétiennes à s'appliquer à orner leur intérieur avec plus de soin que ce qui se montre au-dehors, les exhortant à s'abstenir d'arranger leurs cheveux avec art, de se couvrir d'or, et de mettre trop de recherche dans leur manière de se vêtir\footnote{1. Petr. 3. 4.}. Un des grands reproches faits au mauvais riche dans l'Évangile, est d'avoir porté du luxe et de la somptuosité dans ses habits, et de n'avoir pas fait l'aumône\footnote{Luc. 16. 19, 20, 21.}. Elles sont épouvantables aussi les menaces que Dieu fit autrefois aux filles de Sion par l'organe d'un de ses prophètes, à cause de la fierté de leur démarche, de l'excès de leur parure, du défaut de modération dans leur mise.
 
Instruits à l'école de Jésus-Christ, des prophètes et des prédicateurs de l'Évangile, les Pères de l'Église semblent n'élever qu'une voix contre le luxe des habits, et tout ce qui peut s'y montrer en opposition avec les règles sévères de la plus exacte modestie. Tertullien remarque deux sortes d'excès chez les femmes relativement à notre objet. Il trouve le premier dans l'or, les pierres précieuses, la somptuosité de l'habillement, et il taxe cet excès de crime d'ambition; il découvre l'autre dans l'art recherché qu'elles emploient à leur chevelure, au poli de la peau, et à faire paroître ce qui attire sur elles les regards; et il ne craint pas de donner à ce dernier le nom odieux de prostitution. Selon saint Cyprien, il ne convient qu'aux filles qui ont abjuré toute pudeur, et qu'aux femmes perdues de mœurs, d'étaler des habits et des parures trop ornés, ainsi que les attraits des formes et de la taille du corps. Il ajoute que les personnes qui mettent le plus d'importance à leur mise, sont aussi celles qui font le moins de cas de la vertu. Plus l'homme extérieur, c'est-à-dire le corps, reçoit de culture et d'ornemens, plus l'intérieur est négligé et souffre de dommage, dit saint Augustin. Saint Jérôme assure que les filles qui ornent leur tête, font tomber leurs cheveux sur leur visage, se servent d'artifice pour se donner du poli, se colorent les joues, portent des manches serrées, des robes sans plis, des chaussures ornées de frisure, périssent, en s'affichant comme à l'enchère, tout en conservant le nom de vierges\footnote{Ut, sub nomine virginali vendibilis pereant. S. Hieron.}. Il regarde aussi les jeunes gens qui se frisent et qui se parfument, comme des pestes et des poisons qui tuent la pudeur\footnote{Quasi quasdam pestes et venena pudicitiae. Ibid.}.
 
Mais laissons les citations des Pères qu'il seroit facile d'accumuler ici : les païens eux-mêmes nous donnent des leçons qui ne sont pas moins sévères que celles de ces docteurs de l'Église. Afin de détourner les femmes honnêtes de l'amour du luxe et des vanités, dont les suites sont toujours funestes aux mœurs, l'ancienne Sparte ne permettoit qu'aux courtisanes l'usage des parures ornées d'or et de fleurs. Nous voyons dans un historien grave, qui a écrit sa Vie, que l'empereur Auguste avoit coutume de dire qu'un vêtement où éclatoient la somptuosité et la mollesse, étoit l'étendard de la superbe et le repaire du libertinage\footnote{[illisible] Sueton.}.
 
Le Saint-Esprit nous dit dans l'Écriture que l'habillement, certains ris immodérés, et la démarche de l'homme annoncent ce qu'il est, c'est-à-dire qu'ils le trahissent et qu'ils découvrent l'intérieur de son âme\footnote{Amictus corporis, et risus dentium, et ingressus hominis enuntiant de illo. Eccl. 19. 27.}. Quel jugement peut-on donc porter des personnes qui osent fouler publiquement aux pieds les règles les plus communes de la bienséance et de la modestie naturelle? Ont-elles droit d'exiger qu'on croie à leur vertu, tandis qu'elles produisent au-dehors et avec art ce que la vertu commande de cacher et d'oublier? La recherche de plaire par la beauté, dit encore Tertullien, ne vient point d'une conscience sans tache : on sait que naturellement c'est un piège tendu pour attirer au crime.
 
Et qui ne voit quelle opposition il y a entre une conduite semblable et l'esprit de notre religion sainte? Ignorez-vous, dit saint Paul, que vos corps sont les membres de Jésus-Christ? Prendrai-je donc les membres de Jésus-Christ pour en faire les membres d'une prostituée\footnote{1. Cor. 6. 15.}?
 
Dans le baptême, on a renoncé à Satan, à ses pompes et à ses œuvres : or, n'est-ce pas renouer avec lui, arborer en quelque sorte l'étendard de ses pompes, et travailler de concert avec lui pour perdre les âmes, que de se parer d'une manière si peu décente et si peu conforme à la modestie chrétienne?
 
Serait-il nécessaire de montrer encore combien le luxe et l'immodestie des habits attirent de maux dans les familles? Si une femme a ce luxe en tête, quel ennui ne cause-t-elle point à son époux? quelle dépense peut suffire à son ambition? à quel artifice ne recourt-elle pas pour déguiser, cacher l'emploi qu'elle fait de l'argent qui lui est confié? N'en viendra-t-elle point jusqu'à contracter des dettes, jusqu'à voler dans la bourse du chef, jusqu'à vendre secrètement des effets, des denrées? enfin jusqu'à détourner sur des domestiques ou sur d'autres personnes innocentes, les soupçons odieux qui naîtront tôt ou tard de sa conduite pleine d'artifices et d'infidélités? Mais quelle source féconde de jalousies, de divisions, de troubles et de scandales, s'il arrive que l'époux prenne ombrage d'un luxe si recherché, si condamnable! Parlons-nous de la dureté qu'on a pour les pauvres, auxquels on refuse dédaigneusement l'aumône? de la ruine des ménages que la fureur de vouloir suivre toutes les modes met à la gêne, ou même renverse de fond en comble? Quelle pitié de voir une mère parée, à la tête d'une famille où des enfans à demi nus demandent inutilement du pain, manquent d'instruction et de toute ressource pour en acquérir!
 
Le luxe est aujourd'hui poussé si loin, il est si généralement répandu, qu'on a peine à connoître, à la mise, les différens rangs de la société : la suivante veut briller comme sa maîtresse; la bourgeoise, comme la femme de condition; l'ouvrier, le journalier, le domestique emploie à sa parure tout le produit de son travail, de ses journées, de ses gages.
 
\secundo Cependant on s'étonne et on se plaint des progrès effrayans que font de jour en jour, au milieu de nous, la dépravation des mœurs et l'impiété, qui en est la suite ordinaire. Sur quoi peuvent être fondés cette surprise et cet étonnement? Une des principales causes de ces calamités publiques ne saute-t-elle pas aux yeux de tous ceux qui veulent les ouvrir sur ce qui se passe? des personnes du sexe ne savent plus comment se mettre aujourd'hui pour attirer sur elles les regards, pour faire parade de leur beauté, pour plaire. Elles empruntent du théâtre même leurs modes; elles exposent aux yeux des nudités scandaleuses, ou si elles les déguisent, c'est avec un art plus séduisant encore : il leur faut des robes presque sans plis, qui découvrent leurs épaules, qui laissent paroître leur gorge, qui n'aient presque point de manches, qui dessinent toute la finesse de leur taille : leurs bas mêmes sont à jour, et leur tête est un printemps artificiel. Comment les mœurs pourraient-elles se conserver pures au milieu de tant de pièges, de tant de scandales, de tant de dangers de séduction?
 
Filles et femmes mondaines, qui bravez ainsi toutes les règles de la bienséance et de la modestie, dont l'observation siéroit si bien à votre sexe et en feroit un des plus beaux ornemens, apprenez une fois combien votre vanité et votre manière indécente de vous habiller vous rendent coupables devant Dieu. Vous êtes, dit un ancien docteur, les portes funestes par où le démon entre librement dans les cœurs. Aussi perfides que lui, et non moins criminelles que la première femme, vous tentez les hommes, et vous les portez à des transgressions aussi pernicieuses dans les suites que criminelles en elles-mêmes. Désertant les premières la loi du Seigneur, vous cherchez des complices dans votre révolte; et par vos ruses pleines de perfidie, par le poison attrayant de vos manières séduisantes, le scandale énorme de vos parures immodestes, vous réussissez à tuer la plupart des imprudens qui arrêtent sur vous leurs regards. Oui, elles portent la mort jusqu'au fond du cœur des hommes, ces formes que vous vous plaisez à faire ressortir, et que vous ornez avec tant de soin; vous servez donc de glaive à la mort\footnote{Tu es diaboli janua, tu es arboris illius resignatrix; tu es divinae legis prima desertrix; tu hominem elisisti, et facta es gladius illi. Tert. de hab. mul.}.
 
Qui pourrait calculer le nombre prodigieux de péchés dont un pareil luxe est la source? Qui est-ce qui en portera la peine dans l'éternité malheureuse? sans doute ceux qui ont l'imprudence de les commettre. Mais les misérables qui y auront donné lieu et qui en auront été les premières causes, qu'en sera-t-il d'elles au jugement du souverain juge? Il a déjà prononcé leur sentence : Malheur à l'homme par qui le scandale arrive\footnote{Matth. 8. 7.}!
 
\tertio On ne manque cependant pas de prétextes pour se rassurer. C'est la mode, dit-on; je me couvrirai de ridicule si je m'en écartais trop; je n'ai point d'intention mauvaise; j'en fais moins que beaucoup d'autres qui passent pour honnêtes.
 
De bonne foi, croyez-vous que ces vaines excuses vous exemteront de condamnation au redoutable tribunal de Dieu? C'est la mode; oui, mais cette mode est mauvaise : elle a été inventée au théâtre; elle n'est propre qu'à enflammer des passions toujours prêtes à s'allumer.... Vous encourriez le blâme et vous paroîtriez ridicule, si vous ne vous mettiez pas comme les autres! N'y a-t-il donc plus de modèle de modestie et de retenue dans la ville ou le bourg que vous habitez? Et quand il n'y en auroit point (ce qui seroit bien déplorable), faudrait-il rougir de la vertu, et craindre de reconnoître Jésus-Christ devant les hommes, afin qu'il ne vous reconnoût pas lui-même aux yeux de son Père céleste? De qui seriez-vous blâmée? de quelques personnes légères, imprudentes, que le démon tient dans ses liens, et qui, n'ayant pas le courage de suivre votre exemple, chercheront à se dédommager de ce qui leur porte envie en vous par quelques railleries et quelques sarcasmes. Mais n'est-ce pas ainsi qu'un monde libertin et maudit en agit tous les jours contre quiconque ne partage pas sa perversité, sa corruption? La porte du ciel seroit-elle si étroite, s'il étoit possible d'y entrer en faisant comme la plupart font? Dussiez-vous paroître plus ridicule que vous ne le dites et même que vous ne le pensez, ne devriez-vous pas porter cette ignominie à la suite de Jésus chargé d'humiliations, d'opprobres et d'un gibet infâme, à cause de vous et pour votre amour? Si quelques personnes vous blâment, combien d'autres admireront votre courage, loueront votre vertu, et s'efforceront peut-être de marcher sur vos traces!... Vous n'avez pas de mauvaise intention, cela est-il bien certain? Cherchez jusqu'au fond de votre cœur; portez-y le flambeau de la lumière qui éclaire tout, et qui ne flatte jamais. Voudriez-vous être surprise par la mort dans l'accoutrement où vous osez vous montrer? Ah! qu'il est facile de se faire illusion pendant qu'on jouit de la santé, qu'on est jeune encore, et qu'on se croit très-éloigné du terme de la vie, et du jugement qui suit de près! Mais qu'on pense différemment quand on se voit à la porte de l'éternité, à la veille d'être jugé! Je veux cependant que votre intention ne soit pas mauvaise à vos propres yeux; que même vous ne mettiez aucune complaisance blâmable dans votre parure immodeste, et qu'enfin vous détestiez la tyrannie de la mode, tout en la suivant de trop près : échapperez-vous encore à la rigueur de celui qui juge les justices mêmes? Si vous ne pensez pas mal, vous agissez mal, et ce ne peut être de bonne foi dans une matière où la raison et le sentiment naturel parlent eux-mêmes, quand on ne les a pas encore volontairement étouffés... Enfin, dites-vous, j'en fais moins que beaucoup d'autres qui passent pour honnêtes. Soit : mais vous en faites encore plus que la sévérité de l'Évangile ne vous en permet; vous transgressez donc la règle sur laquelle vous serez jugé : quel sera votre sort?
 
\section{Résolutions}
 
\primo Je m'habillerai désormais selon ma condition; mais un de mes plus grands soins sera de me conformer, dans ma mise, aux règles saintes de la plus exacte modestie. 

\secundo J'aurai constamment de l'horreur pour toute mode immodeste, et je ne m'y conformerai jamais, quoi qu'il puisse m'en coûter. 

\prayer{O mon Dieu! plutôt mourir que de devenir la cause criminelle de la perte d'une seule âme.}
 
\end{document}
